\documentclass[a4paper,12pt]{report}

% =========================================================
% Compilar com XeLaTeX
% =========================================================
\usepackage{fontspec}
\usepackage[portuguese]{babel}
\usepackage{microtype}
\usepackage{graphicx}
\usepackage{hyperref}
\usepackage{amsmath,amssymb}
\usepackage{titlesec}
\usepackage{geometry}
\usepackage{fancyhdr}
\usepackage{parskip}
\usepackage{xcolor}
\usepackage{booktabs}
\usepackage{longtable}
\usepackage{array}
\usepackage{enumitem}
\usepackage{ifthen}
\usepackage{float}
\usepackage{caption}

\IfFontExistsTF{Times New Roman}
  {\setmainfont{Times New Roman}}
  {\setmainfont{TeX Gyre Termes}}

\definecolor{darkteal}{RGB}{0,102,102}

\hypersetup{
    colorlinks=true,
    linkcolor=darkteal,
    filecolor=darkteal,
    urlcolor=darkteal,
    citecolor=darkteal,
    pdftitle={PRAGtico -- Sistema Inteligente de Apoio à Coordenação de Manobras de Pilotagem},
    pdfauthor={André Marques Pereira}
}

\geometry{left=2.5cm,right=2.5cm,top=2.5cm,bottom=2.5cm}

\titleformat{\chapter}[display]
  {\normalfont\huge\bfseries\color{darkteal}}{\chaptername\ \thechapter}{20pt}{\Huge}
\titleformat{\section}
  {\normalfont\Large\bfseries\color{darkteal}}{\thesection}{1em}{}
\titleformat{\subsection}
  {\normalfont\large\bfseries\color{darkteal}}{\thesubsection}{1em}{}

\fancypagestyle{plain}{
  \fancyhf{}
  \fancyfoot[C]{\thepage}
  \renewcommand{\headrulewidth}{0pt}
}
\renewcommand{\headrulewidth}{0pt}
\renewcommand{\footrulewidth}{0pt}
\pagestyle{plain}

\addto\captionsportuguese{%
  \renewcommand{\contentsname}{Índice}
  \renewcommand{\listfigurename}{Lista de Figuras}
  \renewcommand{\listtablename}{Lista de Tabelas}
}

\begin{document}

\begin{titlepage}
    \centering
    \IfFileExists{logoUAb.png}{\includegraphics[width=0.55\textwidth]{logoUAb.png}\par\vspace{1cm}}{}
    {\scshape\LARGE Projeto Final em Engenharia Informática \par}
    \vspace{1.5cm}
    {\huge\bfseries\color{darkteal} PRAGtico -- Sistema Inteligente de Apoio à Coordenação de Manobras de Pilotagem \par}
    \vspace{0.5cm}
    {\Large Capítulo 1 -- Introdução e Especificação Inicial do Projeto \par}
    \vspace{0.5cm}
    {\Large RELATÓRIO \par}
    \vfill
    {\Large\itshape André Marques Pereira \par}
    \vspace{0.5cm}
    {\Large Orientador: Professor José Coelho \par}
    \vspace{0.5cm}
    {\large Março de 2026 \par}
\end{titlepage}

\tableofcontents
\newpage
\listoffigures
\newpage
\listoftables
\newpage

\chapter*{Resumo}
\addcontentsline{toc}{chapter}{Resumo}

O presente projeto propõe o desenvolvimento de um protótipo designado \textit{PRAGtico}, concebido como um sistema inteligente de apoio à coordenação e marcação de manobras de pilotagem em contexto portuário. A proposta assenta num núcleo funcional composto por uma base operacional estruturada, um motor aplicacional de gestão de estados e um assistente inteligente com recuperação de contexto relevante, sendo complementada por uma interface Web para registo, consulta, validação, administração documental e histórico de manobras. O sistema não pretende substituir o piloto nem automatizar a decisão final sobre a manobra; visa antes apoiar a organização da informação, melhorar a rastreabilidade das operações, tornar visíveis conflitos de planeamento e disponibilizar apoio contextualizado à decisão humana. Para esse efeito, combinará dados estruturados do próprio sistema com documentação operacional, regras configuradas, tabelas de marés, dados meteorológicos e outros elementos relevantes para a coordenação. O capítulo apresenta o enquadramento do problema, as necessidades a colmatar, os objetivos, o trabalho proposto, as restrições, os resultados esperados, as tecnologias previstas, os requisitos funcionais preliminares, a checklist operacional, o modelo funcional preliminar, a representação UML inicial do sistema, a arquitetura conceptual preliminar do assistente inteligente, as fontes de informação previstas, o calendário geral do projeto e a estrutura do restante relatório.

\chapter{Introdução}
\label{chap:introducao}

Este capítulo apresenta a especificação inicial do projeto \textit{PRAGtico -- Sistema Inteligente de Apoio à Coordenação de Manobras de Pilotagem}. Começa por enquadrar o problema e delimitar o âmbito do sistema. Seguidamente, identificam-se as necessidades que motivam o projeto, definem-se os objetivos e descreve-se o trabalho proposto. São ainda apresentadas as principais restrições, os resultados esperados, as tecnologias previstas, os requisitos funcionais preliminares, a checklist operacional, o modelo funcional preliminar, a representação UML inicial do sistema, a arquitetura conceptual preliminar do assistente inteligente, as fontes de informação previstas, o calendário geral de execução e a estrutura do restante relatório. Sempre que adequado, o texto é apoiado por tabelas e por diagramas UML, de modo a complementar a descrição textual do sistema e clarificar o seu funcionamento preliminar.

De forma resumida, o projeto consiste na conceção e desenvolvimento de um sistema de apoio à coordenação de manobras portuárias, centrado num motor de gestão operacional e num assistente inteligente de consulta contextualizada. Sobre este núcleo assentará uma interface Web, através da qual agentes, pilotos e administradores interagirão com o sistema. O agente poderá registar navios e criar pedidos de manobra, que permanecerão em estado \textit{pendente de aprovação} até validação do piloto. O piloto poderá aprovar ou cancelar a manobra, sendo obrigatório registar a razão do cancelamento, preservando-se o histórico para efeitos de rastreabilidade. O administrador será responsável pela gestão documental, configuração do sistema, administração de utilizadores e permissões. O assistente inteligente, por sua vez, utilizará os dados estruturados do sistema e fontes complementares para responder a perguntas operacionais relacionadas com condições da manobra, disponibilidade de recursos e critérios preliminares de viabilidade.

\section{Âmbito do projeto e enquadramento}

A coordenação de manobras portuárias constitui uma atividade exigente, na qual importa articular, num horizonte temporal frequentemente curto, fatores como a natureza do movimento pretendido, as características do navio, a disponibilidade de pilotos, rebocadores e meios de apoio, as condições meteorológicas e de maré, a ocupação dos cais e a coexistência de outras operações previstas para o mesmo período. Em contexto real, esta atividade depende tanto de informação formalmente registada como de conhecimento operacional acumulado, o que torna o processo sensível à dispersão da informação, à pressão temporal e à necessidade de manter rastreabilidade.

No problema que motiva este projeto, o coordenador é um piloto destacado segundo um regime de rotação, responsável por gerir a lista de trabalhos entre as 09h00 de um dia e as 09h00 do dia seguinte. Durante parte significativa desse período, assegura serviço de escritório exigente e pode ainda, posteriormente, integrar a escala do rio e embarcar. Esta conjugação de responsabilidades torna particularmente relevante a existência de ferramentas que permitam reunir informação, tornar visíveis conflitos de planeamento, consolidar critérios de verificação e reduzir o esforço de consulta manual.

É neste contexto que surge o \textit{PRAGtico}. O nome remete para a figura do prático ou piloto da barra e procura refletir o caráter pragmático do sistema, orientado para apoio operacional. O âmbito do projeto será, contudo, deliberadamente delimitado. Não se pretende construir um sistema de controlo automático do navio, nem um sistema de navegação autónoma, nem uma plataforma de produção totalmente integrada com infraestruturas reais. O objetivo é desenvolver um protótipo funcional e academicamente defensável, centrado na coordenação humana da manobra, na organização da informação que a suporta e na demonstração de como um mecanismo inteligente de recuperação de contexto pode acrescentar utilidade a esse processo.

\section{Necessidades que o projeto visa colmatar}

A proposta assenta num conjunto de necessidades concretas observadas no domínio de aplicação:

\begin{enumerate}[label=\arabic*.]
    \item \textbf{Consulta rápida de informação operacional.} Em janelas de decisão curtas, o coordenador necessita de confirmar dados sobre movimentos previstos, condições ambientais, restrições locais, disponibilidade de meios e estado de recursos relevantes.
    \item \textbf{Consolidação de informação dispersa.} A informação relevante encontra-se frequentemente distribuída por tabelas, registos, documentação normativa, fontes públicas e comunicação informal.
    \item \textbf{Rastreabilidade das decisões operacionais.} A aprovação, alteração ou cancelamento de uma manobra deve ficar registado, incluindo o respetivo motivo quando aplicável.
    \item \textbf{Verificação preliminar da viabilidade da manobra.} Antes da aprovação, importa confirmar um conjunto de condições operacionais mínimas, como disponibilidade de pilotos, rebocadores, cais, maré, condições meteorológicas e compatibilidade temporal.
    \item \textbf{Antecipação de conflitos de recursos.} A coordenação de múltiplos movimentos em simultâneo exige visibilidade sobre sobreposições de horários, recursos indisponíveis e constrangimentos espaciais.
    \item \textbf{Redução da carga cognitiva.} A gestão de várias manobras em simultâneo beneficia de instrumentos que reduzam o esforço de pesquisa, integrem dados de origens distintas e organizem o estado global da operação.
    \item \textbf{Demonstração de viabilidade técnica.} É necessário demonstrar, num protótipo controlado, que um sistema deste tipo pode ser útil, tecnicamente exequível e potencialmente adaptável a outros contextos portuários.
\end{enumerate}

\section{Objetivos do projeto}

Os objetivos do projeto decorrem diretamente das necessidades identificadas, procurando estabelecer uma relação clara entre o problema observado e a solução proposta.

\subsection{Objetivo geral}

Desenvolver um protótipo de sistema inteligente de apoio à coordenação de manobras de pilotagem, composto por uma base operacional estruturada, um motor de gestão de estados, uma interface Web e um assistente inteligente de consulta contextualizada, capaz de apoiar o piloto/coordenador na programação, validação e acompanhamento de movimentos portuários.

\subsection{Objetivos específicos}

\begin{enumerate}[label=\arabic*.]
    \item Definir o âmbito funcional do protótipo, os seus utilizadores e os principais fluxos de interação.
    \item Conceber uma interface Web simples com autenticação e perfis diferenciados de administrador, agente e piloto.
    \item Permitir ao agente registar navios e criar pedidos de manobra em estado \textit{pendente de aprovação}.
    \item Permitir ao piloto aprovar ou cancelar manobras, assegurando o registo obrigatório do motivo de cancelamento e a manutenção do histórico.
    \item Permitir ao administrador gerir utilizadores, perfis, permissões, documentação operacional e parâmetros de configuração do sistema.
    \item Definir uma checklist operacional preliminar que avalie fatores como calado, luz do dia, disponibilidade de pilotos e rebocadores, ocupação do cais, maré, meteorologia e conflitos temporais.
    \item Estruturar os dados do sistema de forma a suportar consulta inteligente e rastreabilidade operacional.
    \item Implementar um mecanismo de recuperação e apresentação de informação relevante a partir de dados estruturados, documentação operacional e fontes externas.
    \item Avaliar a utilidade do sistema para responder a questões relativas à viabilidade de manobras, previsão de disponibilidade de recursos e adaptação a outros contextos portuários.
    \item Produzir cenários ilustrativos de teste e analisar criticamente a utilidade e limitações do protótipo em ambiente controlado.
\end{enumerate}

\section{Descrição do trabalho a realizar}

Para atingir os objetivos definidos, o projeto será desenvolvido em várias frentes articuladas. Numa primeira fase, será realizada a análise do domínio e o levantamento dos requisitos funcionais e não funcionais. Nesta etapa, procurar-se-á clarificar o ciclo de vida da manobra, os papéis de administrador, agente e piloto, os dados mínimos a registar por navio e por movimento e os critérios operacionais relevantes para a checklist de verificação.

Numa segunda fase, será desenhada a arquitetura do protótipo. A proposta atual aponta para uma estrutura modular com vários blocos principais: interface Web, autenticação e controlo de acessos, backend aplicacional, base de dados operacional, armazenamento documental, mecanismo de pesquisa vetorial e assistente inteligente. O sistema será concebido de modo a manter um núcleo operacional e cognitivo central, sendo a interface Web a camada principal de interação humana com esse núcleo. A base de dados armazenará navios, pedidos de manobra, estados, histórico de cancelamentos, disponibilidade de recursos, perfis de utilizador, permissões, documentação e parâmetros de configuração. O assistente inteligente recorrerá a estes dados e a fontes externas selecionadas para responder a perguntas operacionais.

Numa terceira fase, será implementada a interface Web. O agente deverá poder criar e consultar pedidos de manobra, introduzindo informação relevante do navio e da operação. O piloto deverá poder consultar as manobras pendentes, aprová-las ou cancelá-las. O administrador deverá poder gerir utilizadores, permissões, documentação e parâmetros estruturais do sistema. Uma manobra cancelada deixará de surgir na lista ativa, mas permanecerá registada em histórico, associada ao motivo introduzido pelo piloto.

Numa quarta fase, será desenvolvida a componente de apoio inteligente. O sistema deverá ser capaz de responder a perguntas como: se um navio com determinadas características pode realizar uma manobra nas condições atuais; se a disponibilidade de pilotos, cais e fundeadouros permite suportar a marcação prevista; ou até que ponto a arquitetura é adaptável a outros contextos portuários mediante substituição ou expansão das fontes de informação. O mecanismo de resposta assentará numa lógica de recuperação de contexto relevante, de modo a reduzir geração livre sem suporte factual.

Por fim, o protótipo será testado com cenários ilustrativos. Um exemplo plausível será a criação, por um agente, de uma manobra de entrada de navio no porto de Setúbal. O sistema registará a operação como \textit{pendente de aprovação}, aplicará a checklist preliminar e disponibilizará ao piloto informação sobre recursos, condições ambientais, ocupação do cais e conflitos temporais, permitindo posteriormente aprovar ou cancelar a manobra com registo da respetiva razão.

\section{Restrições do projeto}

A solução proposta deve ser compreendida à luz de várias restrições que condicionam o seu desenho e a sua execução:

\begin{enumerate}[label=\arabic*.]
    \item \textbf{Âmbito académico.} Trata-se de um projeto individual de licenciatura, pelo que a complexidade técnica deve manter-se compatível com o calendário e recursos disponíveis.
    \item \textbf{Prototipagem.} O sistema será desenvolvido como protótipo demonstrador e não como solução imediatamente pronta para exploração operacional em produção.
    \item \textbf{Integração limitada.} Numa fase inicial, não se pressupõe integração total com sistemas reais do porto, podendo recorrer-se a dados simulados, configuráveis ou parcialmente públicos.
    \item \textbf{Dados disponíveis.} Algumas fontes de informação poderão não estar acessíveis, não existir em formato estruturado ou exigir adaptação para uso em protótipo.
    \item \textbf{Responsabilidade decisória.} A decisão final sobre a manobra continuará sempre a pertencer ao piloto/coordenador, sendo o sistema apenas uma ferramenta de apoio.
    \item \textbf{Validação.} A avaliação do sistema incidirá sobre cenários controlados, não equivalendo a validação operacional em ambiente real.
    \item \textbf{Dependência de serviços externos.} A utilização de APIs de IA, meteorologia ou outras fontes externas poderá estar condicionada por disponibilidade, limites de utilização ou alterações futuras do serviço.
\end{enumerate}

\section{Resultados esperados}

No final do projeto, espera-se obter um conjunto coerente de resultados técnicos e académicos. Em primeiro lugar, pretende-se produzir uma especificação clara do problema, dos requisitos do sistema e das fontes de informação utilizadas. Em segundo lugar, espera-se conceber e implementar um protótipo funcional composto por interface Web, mecanismo de autenticação e perfis, base de dados operacional, gestão de estados das manobras, armazenamento documental e assistente inteligente com recuperação de contexto. Em terceiro lugar, prevê-se disponibilizar uma checklist operacional preliminar, suficientemente estruturada para demonstrar a utilidade de uma verificação automática ou semiautomática antes da aprovação da manobra.

Adicionalmente, espera-se produzir cenários de teste representativos, uma análise crítica das limitações do protótipo e um conjunto de recomendações para evolução futura. Entre essas possibilidades de evolução incluem-se a expansão a outros contextos portuários, a integração de fontes adicionais, o refinamento da componente conversacional e, eventualmente, a introdução de interfaces alternativas, sem alteração do núcleo lógico do sistema. Embora o sistema não seja desenvolvido para colocação imediata em produção, a sua execução deverá demonstrar que uma solução deste tipo poderá, em contexto real e após evolução adicional, apoiar a programação, a rastreabilidade e a verificação preliminar de manobras portuárias.

\section{Tecnologias previstas}

Apesar de a seleção tecnológica poder sofrer ajustamentos ao longo do desenvolvimento, importa desde já indicar a orientação técnica geral do protótipo. Essa explicitação reforça a exequibilidade do projeto e ajuda a compreender a decomposição futura do trabalho.

A Tabela~\ref{tab:tecnologias} sintetiza as tecnologias ou categorias tecnológicas previstas para o desenvolvimento inicial do sistema.

\begin{table}[H]
\centering
\caption{Tecnologias previstas para o protótipo}
\label{tab:tecnologias}
\begin{tabular}{>{\raggedright\arraybackslash}p{4.2cm} >{\raggedright\arraybackslash}p{9.8cm}}
\toprule
\textbf{Componente} & \textbf{Tecnologias previstas / orientação} \\
\midrule
Interface Web & Aplicação Web simples, orientada a formulários, listas de manobras, autenticação e páginas de administração \\
Autenticação & Supabase Auth, com perfis diferenciados de administrador, agente e piloto \\
Backend aplicacional & Serviço responsável pela lógica de negócio, gestão de estados, checklist e mediação entre interface, base de dados, documentação e assistente \\
Base de dados operacional & PostgreSQL para armazenamento estruturado de navios, manobras, estados, histórico, recursos, utilizadores e configurações \\
Pesquisa vetorial / RAG & pgvector para indexação semântica de documentação e suporte à recuperação de contexto relevante \\
Assistente inteligente & LLM acedido por API, com Gemini como opção principal e Cerebras como alternativa de teste e comparação \\
Armazenamento documental & Repositório de documentação operacional, normativa e administrativa acessível ao assistente e ao administrador \\
Fontes externas & Tabelas oficiais de marés, versão estruturada em CSV e dados meteorológicos via API \\
\bottomrule
\end{tabular}
\end{table}

\section{Requisitos funcionais preliminares}

Para tornar o projeto mais explícito do ponto de vista da Engenharia Informática, apresenta-se uma primeira enumeração de requisitos funcionais. Estes requisitos são preliminares e poderão ser refinados após validação do orientador e evolução da análise.

\begin{table}[H]
\centering
\caption{Requisitos funcionais preliminares}
\label{tab:reqfunc}
\begin{tabular}{>{\raggedright\arraybackslash}p{2cm} >{\raggedright\arraybackslash}p{12cm}}
\toprule
\textbf{ID} & \textbf{Descrição} \\
\midrule
RF1 & O sistema deve permitir autenticação diferenciada para utilizadores com perfil de administrador, agente e piloto. \\
RF2 & O agente deve poder registar dados básicos do navio relevantes para a manobra. \\
RF3 & O agente deve poder criar pedidos de manobra com data, hora, tipo de movimento e parâmetros operacionais associados. \\
RF4 & Uma manobra criada por agente deve ficar inicialmente em estado \textit{pendente de aprovação}. \\
RF5 & O piloto deve poder consultar a lista de manobras pendentes e respetivos detalhes. \\
RF6 & O piloto deve poder aprovar uma manobra, alterando o seu estado para \textit{aprovada}. \\
RF7 & O piloto deve poder cancelar uma manobra, sendo obrigatório indicar o motivo do cancelamento. \\
RF8 & Uma manobra cancelada deve deixar de surgir na lista ativa, mantendo-se em histórico para rastreabilidade. \\
RF9 & O sistema deve manter registo dos estados e ações relevantes efetuadas sobre cada manobra. \\
RF10 & O sistema deve permitir verificar critérios operacionais preliminares antes da aprovação da manobra. \\
RF11 & O administrador deve poder gerir documentação relevante para consulta pelo assistente e pelos utilizadores autenticados. \\
RF12 & O administrador deve poder gerir utilizadores, perfis e permissões de acesso ao sistema. \\
RF13 & O assistente deve poder consultar o estado das manobras, recursos, histórico e regras configuradas. \\
RF14 & O assistente deve poder responder a perguntas operacionais com base nos dados estruturados, documentação e fontes externas selecionadas. \\
RF15 & O sistema deve permitir relacionar a resposta do assistente com as fontes e dados utilizados sempre que aplicável. \\
\bottomrule
\end{tabular}
\end{table}

\section{Checklist operacional preliminar}

Uma componente diferenciadora do projeto será a incorporação de uma checklist de verificação preliminar da manobra. Esta checklist não substituirá a análise profissional do piloto, mas servirá para organizar critérios mínimos e tornar visíveis conflitos ou insuficiências antes da aprovação.

A Tabela~\ref{tab:checklist} apresenta um primeiro conjunto de critérios a considerar.

\begin{table}[H]
\centering
\caption{Checklist operacional preliminar da manobra}
\label{tab:checklist}
\begin{tabular}{>{\raggedright\arraybackslash}p{4.2cm} >{\raggedright\arraybackslash}p{9.8cm}}
\toprule
\textbf{Critério} & \textbf{Verificação preliminar} \\
\midrule
Calado do navio & Confirmar compatibilidade com o local, a maré e eventuais restrições operacionais. \\
Luz do dia & Verificar se a manobra exige ou recomenda execução em período diurno. \\
Meios de manobra do navio & Considerar informação relevante sobre propulsão, \textit{bow thruster}, \textit{stern thruster} ou outras capacidades. \\
Rebocadores & Confirmar necessidade, quantidade e disponibilidade temporal de rebocadores. \\
Pilotos disponíveis & Verificar se o número de pilotos disponíveis suporta as manobras planeadas no mesmo intervalo. \\
Disponibilidade do cais & Confirmar se o cais pretendido está livre ou se a ocupação permite acomodar os navios com margens de segurança. \\
Folgas laterais no cais & Considerar, por exemplo, a necessidade de manter cerca de 30 m de separação para evitar contacto entre navios. \\
Maré & Verificar se a manobra depende de reponto ou de janela específica de maré. \\
Bordo pedido & Confirmar se existe requisito particular relativo ao bordo de atracação. \\
Meteorologia & Avaliar impacto de vento forte, chuva intensa ou outras condições adversas. \\
Conflitos temporais & Verificar sobreposição com outras manobras, rebocadores ocupados ou indisponibilidade de recursos. \\
\bottomrule
\end{tabular}
\end{table}

\section{Modelo funcional preliminar}

Com vista a clarificar os principais intervenientes, estados e relações entre os módulos, define-se nesta fase um modelo funcional preliminar do protótipo. Este modelo não substitui a especificação detalhada posterior, mas ajuda a estabilizar a visão do sistema e a estabelecer uma base comum para o desenvolvimento.

Em termos simplificados, o administrador interage com o sistema para gerir utilizadores, permissões, parâmetros e documentação operacional. O agente interage com a interface para registar navios e criar manobras. O piloto interage com a mesma interface para consultar operações pendentes, executar a validação e decidir a aprovação ou cancelamento. O backend aplica as regras de negócio e persiste os dados na base operacional. O assistente inteligente consulta essa base, assim como informação documental complementar, para responder a questões dos utilizadores e apoiar a checklist. A articulação entre estes elementos deverá seguir uma arquitetura modular de mediação de contexto, distinguindo explicitamente dados operacionais, documentação, regras e fontes externas. Nesta fase, contudo, essa organização deve ser entendida como orientação arquitetural preliminar e não como requisito técnico totalmente fechado.

A Tabela~\ref{tab:modelofuncional} sintetiza o modelo funcional preliminar do protótipo.

\begin{table}[H]
\centering
\caption{Modelo funcional preliminar do protótipo}
\label{tab:modelofuncional}
\begin{tabular}{>{\raggedright\arraybackslash}p{3.2cm} >{\raggedright\arraybackslash}p{11cm}}
\toprule
\textbf{Elemento} & \textbf{Descrição} \\
\midrule
Administrador & Utilizador responsável por gerir utilizadores, permissões, documentação e parâmetros de configuração do sistema. \\
Agente & Utilizador responsável por introduzir dados do navio e criar pedidos de manobra. \\
Piloto & Utilizador responsável por analisar, aprovar ou cancelar manobras e registar motivo do cancelamento. \\
Interface Web & Camada principal de interação humana com o sistema. \\
Backend aplicacional & Módulo onde reside a lógica de negócio, gestão de estados, checklist e mediação entre dados e assistente. \\
Base de dados operacional & Repositório estruturado para navios, manobras, histórico, recursos, utilizadores e configurações. \\
Armazenamento documental & Repositório de documentação operacional, normativa e administrativa consultável pelo assistente. \\
Assistente inteligente & Componente de consulta que utiliza dados estruturados, documentação e fontes externas para responder a perguntas operacionais. \\
Checklist operacional & Conjunto de verificações preliminares de apoio à aprovação da manobra. \\
\bottomrule
\end{tabular}
\end{table}

\section{Representação UML preliminar do sistema}

Para complementar a descrição textual do projeto, apresentam-se de seguida dois diagramas UML preliminares. O primeiro descreve o fluxo operacional associado ao ciclo de vida de uma manobra no sistema \textit{PRAGtico}, incluindo a perspetiva administrativa necessária à gestão documental e de acessos. O segundo representa a interação entre os principais intervenientes e componentes do sistema, nomeadamente a interface Web, a autenticação e controlo de acessos, o motor de verificação, o assistente inteligente e as fontes de informação utilizadas.

Estes diagramas não esgotam a modelação do sistema, mas permitem clarificar, desde esta fase inicial, tanto a lógica funcional do processo de manobra como a articulação entre os vários elementos do protótipo. A sua utilização no presente capítulo visa, assim, reforçar a especificação inicial do projeto e preparar o trabalho de desenho mais detalhado a desenvolver nos capítulos seguintes.

\subsection{Diagrama de atividades do processo de manobra}

A Figura~\ref{fig:atividade-manobra} apresenta o diagrama de atividades do processo de gestão de uma manobra no \textit{PRAGtico}. O diagrama evidencia o papel do agente, do sistema, da checklist operacional, do piloto e do administrador, desde o registo inicial da manobra até à decisão final de aprovação ou cancelamento, incluindo ainda os fluxos administrativos de gestão documental e de utilizadores.

\begin{figure}[H]
    \centering
    \includegraphics[width=\textwidth,height=0.78\textheight,keepaspectratio]{diagrama_atividade_manobra.png}
    \caption{Diagrama de atividades do processo de gestão de manobra no sistema PRAGtico}
    \label{fig:atividade-manobra}
\end{figure}

Como se observa na Figura~\ref{fig:atividade-manobra}, a manobra é inicialmente registada pelo agente e colocada em estado \textit{pendente de aprovação}. Seguidamente, o sistema executa uma verificação preliminar com base numa checklist operacional, contemplando critérios como calado, disponibilidade de pilotos e rebocadores, condições de maré, meteorologia, disponibilidade do cais e conflitos temporais. Após esta verificação, a manobra é submetida à apreciação do piloto, que poderá aprová-la ou cancelá-la. Em caso de cancelamento, o sistema exige o registo do respetivo motivo e preserva o histórico da operação.

\subsection{Diagrama de sequência da interação entre componentes}

A Figura~\ref{fig:sequencia-sistema} apresenta o diagrama de sequência preliminar do sistema. Este diagrama ilustra a interação entre os utilizadores do sistema e os seus principais componentes, destacando o papel da interface Web como ponto de entrada principal, da autenticação e controlo de acessos como mecanismo de validação de perfis, do motor de checklist como elemento de validação operacional e do assistente inteligente como mecanismo de apoio à consulta contextualizada.

\begin{figure}[H]
    \centering
    \includegraphics[width=\textwidth,height=0.78\textheight,keepaspectratio]{diagrama_sequencia_pragtico.png}
    \caption{Diagrama de sequência da interação entre interface, utilizadores, assistente e fontes de informação}
    \label{fig:sequencia-sistema}
\end{figure}

De acordo com a Figura~\ref{fig:sequencia-sistema}, o agente utiliza a interface para registar navios e criar manobras, que são armazenadas na base de dados e sujeitas a uma verificação preliminar. O piloto, por sua vez, consulta as manobras pendentes e pode interagir com o assistente inteligente para obter esclarecimentos adicionais com base nos dados do sistema, nas regras configuradas e nas fontes externas de informação, como maré e meteorologia. Em paralelo, o administrador autentica-se com perfil próprio e assegura a manutenção dos utilizadores, permissões e documentos de suporte ao sistema. Por fim, o piloto decide aprovar ou cancelar a manobra, sendo essa decisão devidamente refletida no estado da operação e no respetivo histórico.

\subsection{Síntese da modelação apresentada}

Os dois diagramas apresentados são complementares. O diagrama de atividades enfatiza a perspetiva processual do sistema, isto é, a sequência de ações e decisões associadas ao ciclo de vida de uma manobra. Já o diagrama de sequência evidencia a perspetiva estrutural e comunicacional, mostrando como os diferentes intervenientes e componentes colaboram para suportar o funcionamento global do protótipo.

Esta dupla representação contribui para uma melhor compreensão do projeto e reforça a sua coerência conceptual, ao mostrar que o \textit{PRAGtico} não é apenas uma interface de registo de operações nem apenas um assistente conversacional isolado, mas antes um sistema integrado de apoio à programação, verificação e consulta operacional em contexto de pilotagem portuária.

\section{Arquitetura conceptual preliminar do assistente inteligente}

Embora a interface Web constitua a camada principal de operação e administração do sistema, o núcleo diferenciador do projeto reside no assistente inteligente e no mecanismo de recuperação de contexto que articula dados operacionais, documentação e fontes externas. Esta componente permite distinguir o \textit{PRAGtico} de uma aplicação meramente administrativa, conferindo-lhe capacidade de interpretar contexto, integrar informação heterogénea e apoiar o utilizador com respostas contextualizadas.

Em termos conceptuais, o assistente não será tratado como um \textit{chatbot} genérico desligado do domínio, mas como uma camada de apoio baseada em recuperação de contexto relevante. A sua lógica seguirá uma abordagem do tipo \textit{Retrieval-Augmented Generation} (RAG), na qual a resposta é produzida com base em dados efetivamente recuperados a partir de diferentes fontes, em vez de depender apenas de geração livre. Entre essas fontes incluem-se os dados estruturados do sistema, a documentação operacional selecionada, as regras e configurações, a checklist operacional e, quando aplicável, fontes externas como marés e meteorologia.

Adicionalmente, o assistente deverá respeitar um conjunto de regras de controlo e \textit{guardrails}. Estas restrições têm como finalidade limitar a resposta ao contexto efetivamente disponível, distinguir informação confirmada de inferência, impedir que o assistente substitua a decisão formal do piloto e assegurar o respeito pelos perfis e permissões dos utilizadores. Desta forma, o assistente é concebido como mecanismo de apoio à decisão e à consulta, e não como entidade autónoma de validação operacional.

\subsection{Diagrama conceptual do assistente inteligente}

A Figura~\ref{fig:arquitetura-bot} apresenta uma visão conceptual preliminar do funcionamento do assistente inteligente no contexto do sistema \textit{PRAGtico}. O diagrama evidencia que os diferentes perfis de utilizador --- agente, piloto e administrador --- interagem com o sistema através da interface Web, sendo a resposta construída a partir da articulação entre autenticação e controlo de acessos, dados estruturados, documentação, regras operacionais, checklist e fontes externas.

\begin{figure}[H]
    \centering
    \includegraphics[width=\textwidth,height=0.72\textheight,keepaspectratio]{diagrama_arquitetura_bot.png}
    \caption{Arquitetura conceptual preliminar do assistente inteligente do sistema PRAGtico}
    \label{fig:arquitetura-bot}
\end{figure}

Como se observa na Figura~\ref{fig:arquitetura-bot}, a interface Web funciona como ponto de entrada para a consulta do utilizador, encaminhando o pedido para o assistente inteligente e para os restantes componentes do sistema. O assistente, por sua vez, recolhe contexto a partir dos dados estruturados, da documentação operacional, das configurações e regras associadas à checklist e, quando necessário, de fontes externas. Antes da apresentação da resposta final, intervém uma camada de regras e \textit{guardrails}, responsável por validar coerência, limitar extrapolações indevidas e reforçar a fiabilidade do comportamento do sistema.

\subsection{Síntese do papel do assistente no projeto}

A inclusão desta arquitetura conceptual reforça a distinção entre duas dimensões complementares do projeto: por um lado, a dimensão transacional e operacional do sistema, responsável por registar, organizar e manter o estado das manobras; por outro, a dimensão cognitiva do assistente, orientada para interpretação contextual, apoio à consulta e organização da informação relevante para decisão humana.

Assim, o \textit{PRAGtico} deve ser entendido como um sistema integrado que combina persistência operacional, gestão de estados, verificação preliminar, consulta contextualizada e apoio inteligente à coordenação de manobras, sendo a interface Web o principal meio de acesso a esse núcleo funcional.

\section{Fontes de informação previstas}

Uma vez que o protótipo assenta simultaneamente em gestão operacional e em consulta inteligente, a definição das fontes de informação constitui um aspeto central da sua viabilidade. Nesta fase inicial, pretende-se trabalhar com fontes delimitadas, acessíveis e compatíveis com experimentação controlada, evitando dependência excessiva de integrações complexas.

As fontes previstas podem ser organizadas nas seguintes categorias:

\begin{enumerate}[label=\arabic*.]
    \item \textbf{Dados operacionais estruturados do próprio sistema.} Incluem navios, pedidos de manobra, estados, histórico de cancelamentos, recursos, perfis de utilizador, permissões e parâmetros configuráveis.
    \item \textbf{Documentação normativa e operacional.} Incluem regulamentos, procedimentos internos de pilotagem, regras configuradas, documentação administrativa e outros materiais relevantes para consulta.
    \item \textbf{Dados de marés.} Incluem tabelas oficiais do Instituto Hidrográfico e a sua conversão para formato estruturado em CSV, permitindo integração direta na lógica do sistema e em cenários de teste.
    \item \textbf{Dados meteorológicos.} Incluem informação obtida através de API para vento e condições atmosféricas relevantes à avaliação preliminar da manobra.
    \item \textbf{Cenários ilustrativos de teste.} Sempre que necessário, serão produzidos dados exemplificativos e situações controladas, de modo a assegurar demonstrabilidade e avaliação do protótipo.
\end{enumerate}

A Tabela~\ref{tab:fontesinfo} sintetiza as fontes inicialmente previstas.

\begin{table}[H]
\centering
\caption{Fontes de informação previstas}
\label{tab:fontesinfo}
\begin{tabular}{>{\raggedright\arraybackslash}p{4cm} >{\raggedright\arraybackslash}p{9cm}}
\toprule
\textbf{Categoria} & \textbf{Exemplos / utilização prevista} \\
\midrule
Dados estruturados internos & Navios, manobras, estados, histórico, recursos, utilizadores, permissões, configurações \\
Documentação operacional & Procedimentos internos, documentação administrativa, regras configuradas, materiais de suporte \\
Regulamentação & COLREGs e outras referências normativas relevantes para enquadramento documental \\
Marés & Tabelas oficiais do Instituto Hidrográfico e dataset CSV estruturado para consulta automática \\
Meteorologia & Dados de vento e condições atmosféricas via API \\
Cenários de teste & Dados exemplificativos, manobras simuladas e situações controladas para validação \\
\bottomrule
\end{tabular}
\end{table}

\section{Questões orientadoras da investigação}

Para além da implementação técnica do protótipo, o projeto procurará explorar um conjunto de questões orientadoras que ajudam a estruturar a avaliação da solução proposta.

A primeira questão consiste em determinar se um navio com determinadas características pode realizar a manobra nas condições atuais. Esta questão articula dados do navio, marés, meteorologia, recursos disponíveis e outros critérios operacionais, constituindo um exemplo direto de apoio à decisão.

A segunda questão consiste em verificar se o sistema permite prever a disponibilidade de recursos portuários, como pilotos, cais e fundeadouros, para apoio à marcação de manobras. Aqui procura-se avaliar a utilidade do sistema na antecipação de conflitos, na organização temporal das operações e na redução de ambiguidades de coordenação.

A terceira questão consiste em analisar se a solução é adaptável a outros contextos portuários através da atualização das fontes de dados e conhecimento. Esta questão relaciona-se com a modularidade da arquitetura e com a utilização de RAG como mecanismo de especialização do assistente sem necessidade de reconstrução integral da solução.

\section{Calendário geral do projeto}

A execução do projeto será faseada, de modo a permitir consolidar progressivamente requisitos, arquitetura, implementação e avaliação. A Tabela~\ref{tab:calendario} apresenta uma proposta inicial de calendarização, a ajustar em função da evolução do trabalho e das orientações do orientador.

\begin{table}[H]
\centering
\caption{Calendário geral preliminar do projeto}
\label{tab:calendario}
\begin{tabular}{>{\raggedright\arraybackslash}p{7.5cm} >{\raggedright\arraybackslash}p{5cm}}
\toprule
\textbf{Fase / Atividade} & \textbf{Período previsto} \\
\midrule
Definição final do problema, âmbito e requisitos & Março de 2026 \\
Levantamento e organização de dados, regras e fontes de informação & Março -- Abril de 2026 \\
Desenho da arquitetura, modelo de dados e autenticação & Abril de 2026 \\
Implementação da interface Web e gestão de estados das manobras & Abril -- Maio de 2026 \\
Implementação da checklist operacional e regras preliminares & Maio de 2026 \\
Integração da base documental, RAG e assistente inteligente & Maio -- Junho de 2026 \\
Desenvolvimento de cenários de teste e demonstração & Junho de 2026 \\
Avaliação, análise crítica e redação final do relatório & Junho de 2026 \\
\bottomrule
\end{tabular}
\end{table}

\section{Estrutura do restante relatório}

Após este capítulo introdutório, o relatório será organizado da seguinte forma:

\begin{itemize}
    \item \textbf{Capítulo 2 -- Desenho:} apresentação da arquitetura do sistema, do modelo de dados, dos diagramas UML, do mecanismo de autenticação e perfis, das tecnologias escolhidas, da checklist operacional e das principais decisões de desenho.
    \item \textbf{Capítulo 3 -- Implementação:} descrição da decomposição do trabalho em tarefas, estrutura da interface, integração com base de dados, implementação da lógica de negócio, mecanismo RAG e assistente inteligente.
    \item \textbf{Capítulo 4 -- Testes:} definição dos cenários de teste, análise dos resultados obtidos e avaliação do comportamento do protótipo.
    \item \textbf{Capítulo 5 -- Conclusões:} síntese crítica do trabalho realizado, limitações identificadas e propostas de evolução futura.
\end{itemize}

Poderão ainda ser incluídos anexos com código, capturas de ecrã, tabelas de apoio, diagramas complementares e exemplos de cenários de utilização.

\chapter{Desenho}
\section*{Texto provisório}
\addcontentsline{toc}{section}{Texto provisório}
Capítulo a desenvolver em fase posterior, após validação da especificação inicial do projeto.

\chapter{Implementação}
\section*{Texto provisório}
\addcontentsline{toc}{section}{Texto provisório}
Capítulo a desenvolver em fase posterior.

\chapter{Testes}
\section*{Texto provisório}
\addcontentsline{toc}{section}{Texto provisório}
Capítulo a desenvolver em fase posterior.

\chapter{Conclusões}
\section*{Texto provisório}
\addcontentsline{toc}{section}{Texto provisório}
Capítulo a desenvolver em fase posterior.

\begin{thebibliography}{99}
\addcontentsline{toc}{chapter}{Bibliografia}

\bibitem{uabmodelo}
Universidade Aberta. (s.d.). \textit{Modelo de relatório do Projeto Final em Engenharia Informática}. Documento de apoio da unidade curricular.

\bibitem{ihmares}
Instituto Hidrográfico. (2026). \textit{Tabela de Marés 2026}. Disponível em: \url{https://loja.hidrografico.pt/ln/web/wp-content/uploads/2023/11/TabelaMare_I_2026_signed.pdf}. Consultado em março de 2026.

\bibitem{colreg}
Instituto Hidrográfico. (s.d.). \textit{Regulamento Internacional para Evitar Abalroamentos no Mar}. Disponível em: \url{https://loja.hidrografico.pt/sdm_downloads/regulamento-internacional-para-evitar-abalroamentos-no-mar/}. Consultado em março de 2026.

\bibitem{kagglemares}
Pereira, A. (2026). \textit{Portuguese Ports Tides 2024}. Kaggle. Disponível em: \url{https://www.kaggle.com/datasets/andpereira/portuguese-ports-tides-2024}. Consultado em março de 2026.

\bibitem{weatherapi}
WeatherAPI. (s.d.). \textit{Weather API for Weather Data}. Disponível em: \url{https://www.weatherapi.com/}. Consultado em março de 2026.

\bibitem{supabase}
Supabase. (s.d.). \textit{Supabase Documentation}. Disponível em: \url{https://supabase.com/}. Consultado em março de 2026.

\bibitem{gemini}
Google AI for Developers. (s.d.). \textit{Gemini API Documentation}. Disponível em: \url{https://ai.google.dev/}. Consultado em março de 2026.

\bibitem{cerebras}
Cerebras. (s.d.). \textit{Cerebras Cloud}. Disponível em: \url{https://cloud.cerebras.ai/}. Consultado em março de 2026.

\bibitem{rag}
Lewis, P., Perez, E., Piktus, A., Petroni, F., Karpukhin, V., Goyal, N., K{\"u}ttler, H., Lewis, M., Yih, W.-t., Rockt{\"a}schel, T., Riedel, S., \& Kiela, D. (2020). \textit{Retrieval-Augmented Generation for Knowledge-Intensive NLP Tasks}. Advances in Neural Information Processing Systems, 33, 9459--9474.

\bibitem{pgvector}
pgvector. (s.d.). \textit{Open-source vector similarity search for Postgres}. Disponível em: \url{https://github.com/pgvector/pgvector}. Consultado em março de 2026.

\bibitem{dss}
Power, D. J. (2002). \textit{Decision Support Systems: Concepts and Resources for Managers}. Westport: Quorum Books.

\end{thebibliography}

\appendix
\chapter{Anexo I -- Cenários de exemplo}
Espaço reservado para inclusão de cenários de teste, tabelas operacionais e exemplos de perguntas ao sistema.

\chapter{Anexo II -- Elementos técnicos}
Espaço reservado para capturas de ecrã, diagramas, código, estruturas de dados ou outros artefactos do projeto.

\end{document}